% Section 5 - Condition variables
% Roberto Masocco <roberto.masocco@uniroma2.it>
% April 22, 2026

% ### Condition variables ###
\section{Condition variables}
\graphicspath{{figs/section5/}}

% --- Condition variables ---
\begin{frame}{Condition variables}{The waiting problem}
	\justifying
	Sometimes a thread must wait for a specific \textbg{condition} to become true before it can proceed — for example:
	\begin{itemize}
		\item a consumer thread waiting until a \textbg{queue is non-empty};
		\item a worker waiting until a \textbg{task is available};
		\item a thread waiting for a \textbg{configuration to be loaded}.
	\end{itemize}
	\bigskip
	\textbg{Polling} the condition in a busy loop (\textbg{spin-waiting}) wastes CPU cycles.\\
	\textbg{Sleeping} for a fixed time risks missing the event or adding unnecessary latency.\\
	\bigskip
	\begin{block}{Definition of condition variable}
		A \textbf{condition variable} is a synchronization primitive that allows a thread to \textbf{sleep efficiently} until another thread \textbf{signals} that a condition of interest has changed.
	\end{block}
\end{frame}
\begin{frame}{Condition variables}{\texttt{std::condition\_variable}}
	\justifying
	\texttt{std::condition\_variable} is the C++ wrapper around \texttt{pthread\_cond\_t}.\\
	It works together with a \texttt{std::mutex} and a \texttt{std::unique\_lock}:
	\begin{itemize}
		\item \texttt{wait(lock)}: atomically releases the mutex and puts the thread to sleep; when awoken, it re-acquires the mutex before returning;
		\item \texttt{notify\_one()}: wakes up one waiting thread;
		\item \texttt{notify\_all()}: wakes up all waiting threads.
	\end{itemize}
\end{frame}
\begin{frame}[fragile]{Condition variables}
	\begin{columns}
		\column{.9\textwidth}
		\begin{lstlisting}[style=codebeamer, language=C++, caption=Basic condition variable pattern.]
std::mutex mtx;
std::condition_variable cv;
bool ready = false;

// WAITING THREAD
std::unique_lock<std::mutex> lock(mtx);
cv.wait(lock, [&]{ return ready; }); // captures all by ref.

// SIGNALING THREAD
{
  std::lock_guard<std::mutex> lock(mtx);
  ready = true;
}
cv.notify_one();
    \end{lstlisting}
	\end{columns}
\end{frame}
\begin{frame}{Condition variables}{Spurious wakeups}
	\justifying
	A thread waiting on a condition variable may be \textbg{woken up without a corresponding notification} — this is a \textbg{spurious wakeup}, an artifact of the underlying implementation.\\
	\bigskip
	For this reason, the condition must \textbg{always be rechecked after waking up}.\\
	The overload of \texttt{wait()} that accepts a \textbg{predicate} (a callable returning \texttt{bool}) handles this automatically:
	\begin{itemize}
		\item \texttt{cv.wait(lock, predicate)} is equivalent to:\\
			\texttt{while (!predicate()) cv.wait(lock);}
	\end{itemize}
	\bigskip
	\begin{alertblock}{Always use a predicate}
		\justifying
		Calling \texttt{wait(lock)} without a predicate is a common source of subtle bugs: the thread wakes up, assumes the condition is true, and proceeds — but it may not be.
		\textbr{Always pass a predicate.}
	\end{alertblock}
\end{frame}
\begin{frame}{Condition variables}{\texttt{notify\_one()} vs \texttt{notify\_all()}}
	\justifying
	\texttt{notify\_one()} wakes up \textbg{one} waiting thread (chosen by the OS).\\
	Use it when a single event satisfies a single waiter — e.g., one item added to a queue, one consumer to wake.\\
	\bigskip
	\texttt{notify\_all()} wakes up \textbg{all} waiting threads.\\
	Use it when the condition change may be relevant to multiple waiters — e.g., a shutdown signal that all workers must see.\\
	\bigskip
	\textbg{After waking, each thread re-acquires the mutex and re-checks the predicate.}\\
	If the condition is no longer true (another thread consumed the event first), the thread goes back to sleep automatically — the predicate loop handles this.\\
	\bigskip
	\begin{block}{}
		\centering
		When in doubt, prefer \texttt{notify\_all()} — it is safe,\\
    and the performance difference is negligible for most applications.
	\end{block}
\end{frame}
